\section*{Problem 1}

给定数据矩阵 \[
X = \begin{bmatrix}
    1 & 2 & 3 & 4 \\
    2 & 1 & 0 & 3 \\
\end{bmatrix}
\]
\begin{enumerate}
    \item 求两个特征的样本均值，并写出中心化后的数据矩阵 $\tilde{X}$；
    \item 求 $\tilde{X}\tilde{X}^\top$ 的第一主成分方向。
\end{enumerate}

\section*{Problem 2}

设随机向量 $\boldsymbol{x}$ 的协方差矩阵为 \[
\Sigma = \begin{bmatrix}
    5 & 2 \\
    2 & 2 \\
\end{bmatrix}.
\]
\begin{enumerate}
    \item 证明 \[\mathrm{Var} \left(\boldsymbol{a}^\top \boldsymbol{x} \right) = \boldsymbol{a}^\top \Sigma \boldsymbol{a}.\]
    \item 在约束 $\left\|\boldsymbol{a}\right\| = 1$ 下，求使方差最大的方向。
    \item 写出最大方差，并说明其所对应方向的意义。
\end{enumerate}

\section*{Problem 3}

设协方差矩阵 \[\Sigma = \begin{bmatrix}
    4 & 1 & 0 \\
    1 & 4 & 0 \\
    0 & 0 & 1 \\
\end{bmatrix}.\]
\begin{enumerate}
    \item 求 $\Sigma$ 的全部特征值及对应的一组单位特征向量。
    \item 求前 1 个、前 2 个主成分捕获的总方差。
    \item 计算前 1 个、前 2 个主成分的方差解释比例，并说明其意义。
\end{enumerate}
